\documentclass[10pt,a4paper]{amsart}
\usepackage[margin=19mm]{geometry}
\usepackage{amsmath,amssymb,mathtools}
\usepackage[scr=rsfs,cal=euler]{mathalfa}
\usepackage{xfrac}
\usepackage[unicode,hidelinks,bookmarks=false]{hyperref}
\newcommand{\ie}{i.e.\ }
\newcommand{\cf}{cf.\ }
\newcommand{\resp}{respectively\ }
\newcommand{\Subsection}{\subsection}
\providecommand{\nobreakdash}{\nobreak\hbox{-}}
\setlength{\emergencystretch}{2em}
\multlinegap0pt
\usepackage{amssymb}
\usepackage{enumerate}
\usepackage{bm}
\usepackage{xcolor}
\newtheorem{claim}{Claim}
\title{Notes on acceptable bundles II}
\author{Osamu Fujino, Taro Fujisawa, Takashi Ono}
\date{Source: arXiv:2604.05233v1 (CC BY 4.0)}
\begin{document}
\maketitle

\noindent\emph{Source and licence.} Notes on acceptable bundles II. Version arXiv:2604.05233v1 (CC BY 4.0), distributed by arXiv under the Creative Commons Attribution 4.0 International licence, http://creativecommons.org/licenses/by/4.0/, as recorded in the arXiv licence metadata for this exact version and shown on the abstract page https://arxiv.org/abs/2604.05233v1. Full licence notice: \texttt{licenses/arxiv-2604.05233-CC-BY-4.0.txt}.\par\smallskip

\noindent\textit{Editorial scope.} Counted unit: the article's claim x-claim12.2 (source0.tex lines 2799--2803). The article's complete original proof of it is delivered with it (source0.tex lines 2805--2865, one \texttt{proof} environment of 61 lines). No mathematics is written, repaired or re-derived: the statement and the proof are the article's own lines, in the article's own notation, and no retained line is substituted or re-flowed.  No further source-local context is needed: every cross-reference of every retained line resolves inside the two delivered spans.  Nothing was omitted from any retained span, so the package carries no omission record.  No retained line contains a citation, so the wrapper carries no bibliography and no bibliography entry.  No retained line contains a figure environment, an image-inclusion command or drawing code, so no image is delivered and the package declares no asset dependency.  Editorial formatting only, in addition: an AMS \texttt{amsart} preamble that restates the article's own declarations and macros used by the retained lines, namely the environments \texttt{claim} on separate counters, together with the standard packages those lines need.  The retained environments print in this document as Claim 1 in order of appearance, whereas the article prints the same items with the numbering of its own full document.  The complete native source of the article as distributed in the arXiv e-print for this exact version is bundled unchanged as \texttt{sources/197942/source0.tex}.
\begin{claim}\label{x-claim12.2}
For some $0<R<1$, the families $\{V_1,\ldots,V_r\}$ and
$\{W_1,\ldots,W_r\}$ form frames of ${}_\alpha E$ and
${}_{1-\alpha-\varepsilon}(E^\vee)$ on $X(R)$, respectively.
\end{claim}

\begin{proof}[Proof of Claim \ref{x-claim12.2}]
Note that, by construction,
\[
V_i \cdot W_j \in {}_{1-\varepsilon}\left(\mathcal O_{X^*}\right)
= \mathcal O_X,
\qquad
(V_i \cdot W_j)\big|_{\pi^{-1}(P)} = \delta_{ij}.
\]
Let
\[
A := (V_i \cdot W_j)_{i,j}
\]
be the associated $r\times r$ matrix-valued holomorphic function.
Shrinking the radius if necessary, we may assume that
\[
\det A \neq 0 \quad \text{on } X(R)
\]
for some $0<R<1$.
Define
\[
\begin{pmatrix}
W'_1 \\ \vdots \\ W'_r
\end{pmatrix}
=
A^{-1}
\begin{pmatrix}
W_1 \\ \vdots \\ W_r
\end{pmatrix}.
\]
Replacing $W_i$ by $W'_i$, we may further assume that
\[
V_i \cdot W_j = \delta_{ij}
\quad \text{on } X(R).
\]
We also assume that
\[
V_1 \wedge \cdots \wedge V_r \neq 0
\quad \text{on } X^*(R).
\]

Let $\Phi \in {}_\alpha E$.
Then $\Phi$ admits the expansion
\[
\Phi = \sum_{i=1}^r (\Phi \cdot W_i)\, V_i,
\]
where
\[
\Phi \cdot W_i \in {}_{1-\varepsilon}\left(\mathcal O_{X^*}\right)
= \mathcal O_X.
\]
It follows that $\{V_1,\ldots,V_r\}$ forms a local frame of
${}_\alpha E$ on $X(R)$.

Similarly, one checks that $\{W_1,\ldots,W_r\}$ is a local frame of
${}_{1-\alpha-\varepsilon}(E^\vee)$.
In particular, we obtain the identification
\[
{}_{1-\alpha-\varepsilon}(E^\vee) = ({}_\alpha E)^\vee.
\] 
We complete the proof. 
\end{proof}

\end{document}
